\DocumentMetadata{
  pdfversion=2.0,
  pdfstandard=ua-2,
  lang=en-US,
  tagging=on,
  tagging-setup={math/setup=mathml-SE,tabsorder=structure}
}
\documentclass[11pt,letterpaper]{article}
\usepackage[margin=0.68in,includefoot]{geometry}
\usepackage{newcomputermodern}
\usepackage{amsmath,amscd,mathtools}
\providecommand{\square}{\mdlgwhtsquare}
\usepackage[hidelinks]{hyperref}
\hypersetup{
  pdftitle={Analysis PhD Exam, May 2021},
  pdfauthor={University of Florida Department of Mathematics},
  pdfsubject={Graduate mathematics examination},
  pdfdisplaydoctitle=true,
  bookmarksopen=true
}
\setcounter{secnumdepth}{0}
\setlength{\parindent}{0pt}
\setlength{\parskip}{0.28em}
\setlength{\emergencystretch}{3em}
\setlength{\leftmargini}{2.15em}
\setlength{\leftmarginii}{2.65em}
\setlength{\leftmarginiii}{3em}
\renewcommand{\labelenumi}{\textbf{\theenumi.}}
\pagestyle{empty}
\raggedbottom
\allowdisplaybreaks
\newenvironment{examproblems}[1]{%
  \begin{enumerate}%
  \setcounter{enumi}{#1}%
  \setlength{\topsep}{0.35em}%
  \setlength{\partopsep}{0pt}%
  \setlength{\itemsep}{0.48em}%
  \setlength{\parsep}{0.18em}%
}{\end{enumerate}}
\newenvironment{examsubparts}{%
  \begin{enumerate}%
  \setlength{\topsep}{0.18em}%
  \setlength{\partopsep}{0pt}%
  \setlength{\itemsep}{0.16em}%
  \setlength{\parsep}{0.1em}%
}{\end{enumerate}}
\newenvironment{examinstructions}{%
  \par\begingroup\itshape
}{%
  \par\endgroup\vspace{0.18em}
}
\makeatletter
\renewcommand\section{\@startsection{section}{1}{0pt}%
  {-1.25ex plus -.25ex minus -.1ex}%
  {0.55ex plus .1ex}%
  {\normalfont\large\bfseries}}
\renewcommand{\@maketitle}{%
  \newpage\null\vskip -1.2em%
  {\footnotesize\bfseries
    \href{https://www.ufl.edu/}{University of Florida}, \href{https://math.ufl.edu/}{Department of Mathematics}\par}%
  \vskip 0.18em%
  \tagstructbegin{tag=Title}%
    {\LARGE\bfseries\href{https://gma.math.ufl.edu/exams/analysis-phd/}{Analysis PhD Exam}, May 2021\par}%
  \tagstructend%
  \vskip 0.35em%
  \tagmcbegin{artifact}%
    \hrule height 1pt%
  \tagmcend%
  \vskip 0.55em%
}
\makeatother
\title{Analysis PhD Exam, May 2021}
\author{}
\date{}
\begin{document}
\maketitle
\thispagestyle{empty}
\begin{examinstructions}
Write solutions in a neat and logical fashion, giving complete reasons for all steps. Attempt SIX problems.
\end{examinstructions}
\begin{examproblems}{0}
\item
Let \(f\in L^1(\mathbb{R})\) and suppose that 
\[
\int_a^b f\,dm=0
\]
 for all rational \(a<b\). Does this imply that \(f=0\) a.e.? Prove, or give a counterexample.
\item
Prove that if \(E\subset\mathbb{R}\) has positive Lebesgue measure, then the set 
\[
E-E=\{x-y\mid x,y\in E\}
\]
 contains an open interval centered at~\(0\).
\item
Let~\(\mu\) and~\(\nu\) be nonzero, \(\sigma\)-finite measures defined on the same measurable space \((X,\mathcal{M})\), and suppose each is absolutely continuous with respect to the other. Prove that 
\[
\frac{d\nu}{d\mu}=\frac{1}{\left(\frac{d\mu}{d\nu}\right)}\quad \mu\text{-a.e.}
\]
\item
Let~\(\mu\) be a positive measure and suppose that \(1<p<r<q<\infty\). Prove that 
\[
L^p(\mu)\cap L^q(\mu)\subset L^r(\mu).
\]
\item
Let~\(H\) be a Hilbert space and \(M\subset H\) a norm-closed subspace. Prove that if \((x_n)\) is a sequence from~\(M\) and \(x_n\) converges weakly to \(x\in H\), then in fact \(x\in M\).
\item
This problem has two parts.
\begin{examsubparts}
\item[\textnormal{a)}]
Prove that for each integer \(n\ge 0\), there exists a finite Borel measure \(\mu_n\) on \([0,1]\) so that for every polynomial~\(p\) of degree at most~\(n\), 
\[
p'(0)=\int_0^1 p(x)\,d\mu_n(x).
\]
\item[\textnormal{b)}]
Prove that there is no finite Borel measure~\(\mu\) on \([0,1]\) such that 
\[
p'(0)=\int_0^1 p(x)\,d\mu(x).
\]
 for all polynomials~\(p\).
\end{examsubparts}
\item
This problem has two parts.
\begin{examsubparts}
\item[\textnormal{a)}]
State the Banach Isomorphism Theorem and the Closed Graph Theorem.
\item[\textnormal{b)}]
Using the Banach Isomorphism Theorem (or otherwise), prove the Closed Graph Theorem.
\end{examsubparts}
\item
(Short answer, you do not need to give a detailed proof, just a brief explanation.)
\begin{examsubparts}
\item[\textnormal{a)}]
Give an example, if possible, of a sequence \((f_n)\subset L^2(\mathbb{R})\) such that \(f_n\to 0\) weakly, but not in measure.
\item[\textnormal{b)}]
Give an example to show that the~\(\sigma\)-finiteness hypothesis is necessary in Tonelli’s theorem.
\item[\textnormal{c)}]
If \((X,\mathcal{M})\) is a measurable space, \(\mu\) is a signed measure on \((X,\mathcal{M})\), and~\(\mu\) omits the value \(+\infty\), does~\(\mu\) attain a maximum value on~\(\mathcal{M}\)?
\end{examsubparts}
\end{examproblems}
\end{document}
